\documentclass[UTF8,11pt]{ctexart}
\usepackage[a4paper,margin=25mm]{geometry}
\usepackage{amsmath,amssymb,amsthm,mathtools}
\usepackage[all]{xy}
\usepackage{xurl}
\usepackage{hyperref}
\usepackage{enumitem}
\hypersetup{hidelinks,breaklinks=true}
\setlength{\emergencystretch}{3em}
\newtheorem*{theorem}{Theorem}
\newtheorem*{lemma}{Lemma}
\newtheorem*{proposition}{Proposition}
\newtheorem*{corollary}{Corollary}
\newtheorem*{definition}{Definition}
\newcommand{\Spec}{\operatorname{Spec}}
\newcommand{\Hom}{\operatorname{Hom}}
\newcommand{\End}{\operatorname{End}}
\newcommand{\Gal}{\operatorname{Gal}}
\newcommand{\Cl}{\operatorname{Cl}}
\newcommand{\vol}{\operatorname{vol}}
\newcommand{\covol}{\operatorname{covol}}
\newcommand{\tr}{\operatorname{tr}}
\newcommand{\im}{\operatorname{im}}
\newcommand{\id}{\operatorname{id}}
\newcommand{\coker}{\operatorname{coker}}
\newcommand{\stacksref}[2]{\href{https://stacks.math.columbia.edu/tag/#1}{#2 [#1]}}
\begin{document}
\section*{003\quad 有限域椭圆曲线的Hasse界}
\subsection*{来源与整理范围}
Andrew Sutherland, \emph{18.783 Elliptic Curves}, Spring 2021, Lecture 7: \emph{Point counting}, dated 03/10/2021. MIT OpenCourseWare. 收录第1--2页的 Lemma 7.1、Remark 7.2 与 Theorem 7.3，直至 Hasse 定理证明结束。

实际访问原文：\url{https://ocw.mit.edu/courses/18-783-elliptic-curves-spring-2021/5e4e3bd15c9a81db2ac186628da095bf_MIT18_783S21_notes7.pdf}。
获取日期：2026-09-28。采用 OCW 2021 版，而非种子题单中的 2023 版。

\subsection*{作者原命题与人类证明}
\paragraph{7.1 Separable and inseparable endomorphisms.}
Recall that the Frobenius endomorphism $\pi_E$ is inseparable. In order to prove Hasse's theorem we will need to use the fact that $\pi_E-1$ is separable. This follows from a much more general result: adding a separable isogeny to an inseparable isogeny always yields a separable isogeny. Note that the sum of two separable isogenies need not be separable: in characteristic $p>0$, if we have $a+b=p$ and both $a$ and $b$ prime to $p$, then $[a]$ and $[b]$ are both separable but $[a]+[b]=[a+b]=[p]$ is inseparable.

\begin{lemma}[7.1]
Let $\alpha$ and $\beta$ be isogenies from $E_1$ to $E_2$, with $\alpha$ inseparable. Then $\alpha+\beta$ is inseparable if and only if $\beta$ is inseparable.
\end{lemma}
\begin{proof}
If $\beta$ is inseparable then by Corollary 5.4 we can write $\alpha=\alpha_{\mathrm{sep}}\circ\pi^m$ and $\beta=\beta_{\mathrm{sep}}\circ\pi^n$, where $\pi$ is the $p$-power Frobenius map and $m,n>0$. We then have
\[
\alpha+\beta=\alpha_{\mathrm{sep}}\circ\pi^m+\beta_{\mathrm{sep}}\circ\pi^n
=(\alpha_{\mathrm{sep}}\circ\pi^{m-1}+\beta_{\mathrm{sep}}\circ\pi^{n-1})\circ\pi,
\]
which is inseparable (any composition involving an inseparable isogeny is inseparable because inseparable degrees multiply). If $\alpha+\beta$ is inseparable, then so is $-(\alpha+\beta)$, and $\alpha-(\alpha+\beta)=\beta$ is a sum of inseparable isogenies, which we have just shown is inseparable.
\end{proof}

\paragraph{Remark 7.2.}
Since the composition of an inseparable isogeny with any isogeny is always inseparable, Lemma 7.1 implies that the inseparable endomorphisms in $\End(E)$ form an ideal (provided we view $0$ as inseparable, which we do).

\paragraph{7.2 Hasse's Theorem.}
We are now ready to prove Hasse's theorem.
\begin{theorem}[7.3; Hasse]
Let $E/\mathbb F_q$ be an elliptic curve over a finite field. Then
\[
\#E(\mathbb F_q)=q+1-t,
\]
where $t:=\tr\pi_E$ is the trace of the Frobenius endomorphism $\pi_E$ and $|t|\le2\sqrt q$.
\end{theorem}
\begin{proof}
Recall that we defined $\mathbb F_q$ as the splitting field of $x^q-x$ over $\mathbb F_p$, where $p=\operatorname{char}(\mathbb F_q)$, thus
\[
\mathbb F_q=\{\alpha\in\overline{\mathbb F}_p:\alpha^q-\alpha=0\}
=\{\alpha\in\overline{\mathbb F}_q:\alpha^q-\alpha=0\}
\]
is precisely the subfield of $\overline{\mathbb F}_q$ fixed by the $q$-power Frobenius automorphism $x\mapsto x^q$. The Frobenius endomorphism $\pi_E:E\to E$ is defined by $\pi_E(x:y:z)=(x^q:y^q:z^q)$, therefore
\[
\begin{aligned}
E(\mathbb F_q)&=\{P\in E(\overline{\mathbb F}_q):\pi_E(P)=P\}\\
&=\{P\in E(\overline{\mathbb F}_q):\pi_E(P)-P=0\}=\ker(\pi_E-1),
\end{aligned}
\]
where $1$ denotes the multiplication-by-1 map $[1]\in\End(E)$. The Frobenius endomorphism $\pi_E$ is inseparable and $-1$ is separable, so by Lemma 7.1 the endomorphism $\pi_E-1$ is separable, thus the cardinality of its kernel is equal to its degree (by Theorem 5.8). Therefore
\[
\begin{aligned}
\#E(\mathbb F_q)&=\#\ker(\pi_E-1)=\deg(\pi_E-1)\\
&=(\widehat{\pi_E-1})(\pi_E-1)
=\widehat\pi_E\pi_E+1-(\widehat\pi_E+\pi_E)=q+1-t.
\end{aligned}
\]
It remains only to show that $|t|\le2\sqrt q$.

Consider the endomorphism $r\pi_E-s$ for $r,s\in\mathbb Z$ with $s\ne0$. We have
\[
\begin{aligned}
\deg(r\pi_E-s)&=(\widehat{r\pi_E-s})(r\pi_E-s)
=(\widehat\pi_E\widehat r-\widehat s)(r\pi_E-s)
=(\widehat\pi_Er-s)(r\pi_E-s)\\
&=\widehat\pi_Er^2\pi_E-\widehat\pi_Ers-sr\pi+s^2
=r^2\widehat\pi_E\pi_E-rs(\widehat\pi_E+\pi_E)+s^2\\
&=r^2\deg\pi_E-rs\tr\pi_E+s^2\\
&=r^2q-rst+s^2,
\end{aligned}
\]
where we have used Lemmas 6.11 and 6.12, and the fact that $\mathbb Z$ is in the center of $\End(E)$. Dividing by $s^2$ and noting that $\deg(r-\pi_Es)\ge0$ yields the inequality
\[
q(r/s)^2-t(r/s)+1\ge0,
\]
valid for all rational numbers $r/s$. Now $\mathbb Q$ is dense in $\mathbb R$, so we must have $qx^2-tx+1\ge0$ for all real numbers $x$. It follows that the discriminant $t^2-4q$ cannot be positive, which yields the desired bound $|t|\le2\sqrt q$.
\end{proof}


\subsection*{整理说明（非作者正文）}
本题对应种子题单第8题，改用作者的原命题（任意有限域上的椭圆曲线），不另造一个新证明。$\widehat\alpha$ 表示对偶同源；原文把整数与相应的倍乘自同态识别，因而其度数等式中整数与自同态按这一约定书写。

先修为同源的可分/纯不可分分解（Corollary 5.4）、可分同源核大小等于次数（Theorem 5.8）、对偶同源的加法/复合性质及 $\widehat\alpha\alpha=[\deg\alpha]$（Lecture 6 的相关结果）。这些是明确的标准先修，而非把 Hasse 界或等价的迹界作为先修。核的识别、可分性说明、度数二次式和实二次多项式判别式步骤均完整收录。

原 PDF 第2页在展开式中有一次写 $\pi$ 而非 $\pi_E$；随后非负性行写 $\deg(r-\pi_Es)$，而前面的展开式对象为 $r\pi_E-s$。此处均忠实保留原页字样，不把自行修订伪装成作者原文。仅调整长公式的断行。第1--2页均实际查看原页图像；原 PDF 未能下载到本地，随包仅留转录摘录。

\subsection*{可修改的简短标注（非作者正文）}
\textbf{$c$：}有限域椭圆曲线的点数误差。

\textbf{$a$：}Frobenius 自同态；有理点与自同态核的对应；可分同源；对偶同源；自同态的度数与迹；非负二次型；有理数稠密性。

\textbf{cross\_theory\_status = yes。}Theorem 7.3 的证明先把有限域点数转为 $\pi_E-1$ 的核大小，再转为同源度数；通过自同态环中的度数二次式与实多项式非负性得到点数误差界。位置：原文第1页下半至第2页上半。

标注为可修改、非穷尽初稿。选题保留核心转换与完整推导，不将现成点数界直接应用计作本题证明。
\subsection*{许可与修改记录}
材料来自 MIT OpenCourseWare；原作者与课程见上文。按该站所示 \href{https://creativecommons.org/licenses/by-nc-sa/4.0/}{CC BY-NC-SA 4.0} 许可复用。本转录及其附加整理采用同一许可；不得据此声称作者或 MIT 认可本次整理。2026-09-28：ChatGPT 转录为 TeX，加入题号、来源、整理说明和初稿标注；数学正文保持作者语言及顺序。

\end{document}
